\documentclass[10pt,a4paper]{article}

\usepackage[a4paper,top=20mm,bottom=21mm,left=20mm,right=20mm,headheight=20pt]{geometry}
\usepackage[T1]{fontenc}
\usepackage{newtxtext,newtxmath}
\usepackage{amsmath,bm}
\usepackage{booktabs,tabularx,longtable,array,multirow}
\usepackage{float}
\usepackage[table]{xcolor}
\usepackage{enumitem}
\usepackage{fancyhdr}
\usepackage{microtype}
\usepackage{caption}
\usepackage{hyperref}
\usepackage{titlesec}
\usepackage{pdflscape}
\usepackage{needspace}

\definecolor{navy}{HTML}{17324D}
\definecolor{bluegray}{HTML}{EAF0F5}
\definecolor{rulegray}{HTML}{C9D3DD}
\definecolor{linkblue}{HTML}{245A8D}

\hypersetup{
  colorlinks=true,
  linkcolor=linkblue,
  citecolor=linkblue,
  urlcolor=linkblue,
  hypertexnames=false,
  pdfauthor={},
  pdftitle={Supplementary Material for How Much Fluctuation Is Dangerous? Predicting Tipping Points via Statistically Conditioned Digital Twins}
}

\setlength{\parindent}{0pt}
\setlength{\parskip}{5pt}
\setlength{\tabcolsep}{4.5pt}
\renewcommand{\arraystretch}{1.16}
\setlist[itemize]{leftmargin=5mm,itemsep=2pt,topsep=2pt}
\setlist[enumerate]{leftmargin=6mm,itemsep=2pt,topsep=2pt}

\newcolumntype{Y}{>{\raggedright\arraybackslash}X}
\newcolumntype{Z}{>{\centering\arraybackslash}X}
\newcolumntype{L}[1]{>{\raggedright\arraybackslash}p{#1}}
\newcolumntype{C}[1]{>{\centering\arraybackslash}p{#1}}
\newcommand{\scdt}{\mathrm{SCDT}}
\newcommand{\gt}{\mathrm{GT}}
\newcommand{\rc}{\mathrm{RC}}

\titleformat{\section}{\Large\bfseries\color{navy}}{S\thesection}{0.65em}{}
\titleformat{\subsection}{\large\bfseries\color{navy}}{S\thesubsection}{0.65em}{}
\titleformat{\subsubsection}{\normalsize\bfseries}{\thesubsubsection}{0.65em}{}
\titlespacing*{\section}{0pt}{13pt}{5pt}
\titlespacing*{\subsection}{0pt}{10pt}{3pt}
\titlespacing*{\subsubsection}{0pt}{8pt}{2pt}

\pagestyle{fancy}
\fancyhf{}
\lhead{\small Supplementary Material: Statistically Conditioned Digital Twins}
\rhead{\small \thepage}
\renewcommand{\headrulewidth}{0.4pt}
\renewcommand{\headrule}{\hbox to\headwidth{\color{rulegray}\leaders\hrule height \headrulewidth\hfill}}

\captionsetup{font=small,labelfont=bf,labelsep=period,hypcap=false}

\begin{document}

\pagenumbering{gobble}

\begin{titlepage}
  \thispagestyle{empty}
  \centering
  \vspace*{0.26\textheight}
  {\color{navy}\bfseries\fontsize{24}{31}\selectfont
  Supplementary Material for\par}
  \vspace{9mm}
  {\bfseries\fontsize{20}{27}\selectfont
  ``How Much Fluctuation Is Dangerous?\par
  Predicting Tipping Points via Statistically Conditioned Digital Twins''\par}
  \vfill
  {\large Detailed model equations, statistical coordinates, training protocols,\par
  reservoir settings, outcome criteria, and numerical results\par}
  \vspace*{0.12\textheight}
\end{titlepage}

\pagenumbering{roman}
\tableofcontents
\clearpage
\pagenumbering{arabic}

\section{Scope and notation}

This Supplementary Material records the numerical protocol used for the final
food-chain, voltage-collapse, and Kuramoto experiments. It complements the main
text by specifying the target equations, conditioning statistics, data windows,
reservoir implementation, test coordinates, outcome rules, and numerical
thresholds. The physical bifurcation parameter is used to generate and validate
ground-truth trajectories, but it is not supplied to the SCDT during training or
autonomous prediction.

For a sampled trajectory, $n$ denotes the discrete sample index and $\Delta t$
the sampling interval. Three different initialization intervals must be kept
distinct:

\begin{itemize}
  \item \textbf{Target transient:} the target ODE is integrated before the
  retained record begins. These samples are discarded and are never supplied to
  the reservoir.
  \item \textbf{Training washout:} retained target samples are supplied to the
  reservoir, but the early reservoir states are excluded from the readout
  regression so that dependence on the zero reservoir state is reduced.
  \item \textbf{Prediction warm-up (teacher forcing):} a target segment is
  supplied immediately before autonomous prediction. These samples initialize
  the reservoir state; the subsequent SCDT output is fed back as its next input.
\end{itemize}

The labels \emph{inter}, \emph{pre}, and \emph{post} denote, respectively, a
test point inside the training range, a test point beyond the training range but
before the physical transition, and a point beyond the transition.

\section{Statistically conditioned reservoir}

\subsection{Reservoir update}

Let $\bm{u}_n\in\mathbb{R}^{D}$ be the dynamical input, $\bm{r}_n\in
\mathbb{R}^{N_r}$ the reservoir state, and $\bm{s}\in\mathbb{R}^{2}$ a pair of
statistics that remains constant over one trajectory segment. After the
system-specific normalization $\mathcal{T}_s$, the reservoir evolves as
\begin{equation}
  \bm{r}_{n+1}
  = (1-\alpha)\bm{r}_{n}
  + \alpha\tanh\!\left(
      A\bm{r}_{n}+W_{\mathrm{in}}\bm{u}_{n}
      +W_{\mathrm{stat}}\mathcal{T}_s(\bm{s})
    \right),
  \label{eq:scdt-update}
\end{equation}
where $A$ and the input matrices are fixed after random initialization,
$\alpha$ is the leakage rate, and only the output matrix is trained. The
statistics channel therefore acts as a trajectory-level bias that indexes a
family of reservoir dynamics.

For one-step training, the output weights solve
\begin{equation}
  W_{\mathrm{out}}
  = \arg\min_W
  \sum_{m=1}^{M}\sum_{n\in\mathcal{I}_m}
  \left\|W\bm{z}^{(m)}_{n+1}-\bm{u}^{(m)}_{n+1}\right\|_2^2
  +\beta\|W\|_F^2,
  \label{eq:ridge}
\end{equation}
where $M=5$ is the number of training conditions, $\mathcal{I}_m$ excludes the
washout interval, $\beta$ is the ridge coefficient, and $\bm{z}$ is the
system-specific readout feature in Table~\ref{tab:readout}. During autonomous
prediction, $\widehat{\bm{u}}_{n+1}=W_{\mathrm{out}}\bm{z}_{n+1}$ is fed back
as $\bm{u}_{n+1}$.

\begin{table}[ht]
\centering
\caption{System-specific dynamical variables, statistic normalization, and
readout feature. These details refine the generic readout equation in the main
text.}
\label{tab:readout}
\small
\begin{tabularx}{\textwidth}{L{27mm} L{34mm} Y Y}
\toprule
System & Dynamic input/output & Statistic transform $\mathcal{T}_s$ & Readout feature $\bm{z}$ \\
\midrule
Food chain & $[R,C,P]^{\mathsf T}$
& Training-set z-score, followed by elementwise gain $(0.1,0.1)$ and bias
$(-1.63179994,0.05175686)$
& $\bm{z}=\phi(\bm r)$, where zero-based odd entries are squared:
$\phi_j(r_j)=r_j^2$ for $j=1,3,5,\ldots$, and $r_j$ otherwise. \\
Voltage collapse & $[\delta_m,\omega,\delta,V]^{\mathsf T}$
& Coordinatewise min-max scaling with frozen bounds from the original
10,000-sample training records; the scaled
pair is multiplied through $W_{\mathrm{stat}}$, whose standard deviation is
set by the statistic-scaling hyperparameter
& $\bm{z}=\bm r$. \\
Kuramoto network & Scalar order parameter $r$
& Coordinatewise z-score over the five training conditions
& $\bm{z}=[\bm r^{\mathsf T},\widetilde r,1]^{\mathsf T}$, where
$\widetilde r$ is the standardized current order parameter. \\
\bottomrule
\end{tabularx}
\end{table}

\subsection{Conditioning path and test queries}

Each training trajectory provides a measured pair
$\bm{s}^{(m)}=(s_1^{(m)},s_2^{(m)})$. A least-squares line
\begin{equation}
  s_2=a s_1+b
  \label{eq:conditioning-path}
\end{equation}
describes the training trend. The reported test and rate-sweep queries for all
three systems use their corresponding measured-training fits. The physical
parameter is never an input to the reservoir.

Training coordinates in Tables~\ref{tab:food-coordinates}--\ref{tab:kura-coordinates}
are direct measurements from retained target trajectories. The reported test
coordinates are the actual pairs supplied to the final SCDT. For Food and
Voltage, the first test coordinate is obtained from interpolation, direct
precritical measurement, or parameter-to-statistic extrapolation as specified
below. The second coordinate follows the system-specific fit described below.
The Kuramoto test pairs are treated in the same way after selecting asynchronous
inter/pre trajectories.

\section{Conditioning statistics}

\subsection{Standard deviation}

For a scalar record $x_1,\ldots,x_T$, the population standard deviation used in
the code is
\begin{equation}
  \sigma_x=\sqrt{\frac{1}{T}\sum_{n=1}^{T}(x_n-\bar{x})^2},
  \qquad
  \bar{x}=\frac{1}{T}\sum_{n=1}^{T}x_n.
  \label{eq:std}
\end{equation}
The food-chain coordinate is $\sigma_P$ and the voltage coordinate is
$\sigma_V$.

\subsection{Lag-1 autocorrelation}

The voltage experiment uses
\begin{equation}
  \rho_1(V)=\frac{c_1}{c_0},\qquad
  c_0=\frac{1}{T}\sum_{n=1}^{T}(V_n-\bar V)^2,
  \qquad
  c_1=\frac{1}{T-1}\sum_{n=1}^{T-1}(V_n-\bar V)(V_{n+1}-\bar V).
  \label{eq:ac1}
\end{equation}

\subsection{Permutation entropy}

The food-chain experiment uses normalized permutation entropy of $P(t)$ with
embedding order $d=3$ and delay $\tau=1$. If $p(\pi)$ is the empirical
probability of ordinal pattern $\pi$ among the $d!$ possible patterns, then
\begin{equation}
  H_{\mathrm{perm}}
  =-\frac{1}{\log(d!)}\sum_{\pi}p(\pi)\log p(\pi),
  \qquad d=3,\quad \tau=1.
  \label{eq:permentropy}
\end{equation}
The implementation uses the normalized estimator in the \texttt{antropy}
package~\cite{Bandt2002}.

\subsection{Kuramoto order parameter and topology-aware correlations}

Macroscopic synchronization is measured by
\begin{equation}
  r(t)=\left|\frac{1}{N}\sum_{j=1}^{N}e^{\mathrm{i}\theta_j(t)}\right|.
  \label{eq:order-parameter}
\end{equation}
The SCDT dynamically predicts only the scalar $r(t)$. Node phases are used to
construct the two static conditioning coordinates. Set
$x_i(t)=\cos\theta_i(t)$ and let $\rho_{ij}$ be the Pearson correlation between
$x_i$ and $x_j$ over the retained measurement segment. The code uses
$|\rho_{ij}|$. Nodes at or above the 90th percentile of degree form the hub set
$\mathcal H$, and the remaining nodes form $\mathcal L$. The statistics are
\begin{align}
  S_{\mathrm{int}}
  &=\frac{1}{|\mathcal H|(|\mathcal H|-1)}
    \sum_{\substack{i,j\in\mathcal H\\i\ne j}}|\rho_{ij}|,
  \label{eq:sint}\\
  S_{\mathrm{out}}
  &=\frac{1}{|\mathcal H||\mathcal L|}
    \sum_{i\in\mathcal H}\sum_{j\in\mathcal L}|\rho_{ij}|.
  \label{eq:sout}
\end{align}
These are experiment-specific, topology-aware regime coordinates motivated by
correlation-based functional connectivity and hub synchronization studies
\cite{Ziaeemehr2021,Pereira2010,Vlasov2017}; they are not assumed to be
universal early-warning indicators.

\section{Target dynamical systems}

\subsection{Tri-trophic food chain}

The state is $[R,C,P]^{\mathsf T}$: resource, consumer, and predator. The exact
equations implemented in the final experiment are
\begin{align}
  \dot R&=R\left(1-\frac{R}{K}\right)
  -\frac{x_c y_c C R}{R+R_0}, \\
  \dot C&=x_c C\left(\frac{y_cR}{R+R_0}-1\right)
  -\frac{x_p y_p P C}{C+C_0}, \\
  \dot P&=x_p P\left(\frac{y_pC}{C+C_0}-1\right).
  \label{eq:food}
\end{align}
The parameters are
\begin{equation}
  x_c=0.4,\quad y_c=2.009,\quad x_p=0.08,\quad y_p=2.876,
  \quad R_0=0.16129,\quad C_0=0.5.
\end{equation}
The carrying capacity $K$ is varied only to generate and validate data. The
model is a standard chaotic tri-trophic food-chain benchmark
\cite{McCann1994,Kong2021}.

Target trajectories are integrated by RK45. The solver output is evaluated at
$0.1$-time-unit intervals and every tenth value is retained, giving
$\Delta t=1$. For each candidate trajectory, 2,000 retained samples are
discarded as target transient and 6,000 subsequent samples are retained. The
initial condition is sampled independently as
\begin{equation}
  R(0)\sim U(0.6,1.0),\qquad
  C(0)\sim U(0.15,0.55),\qquad
  P(0)\sim U(0.6,0.7).
\end{equation}
Up to 30 draws are attempted. A retained training record is accepted only when
$\min P>0.01$ and its final value satisfies $P>0.01$.

\subsection{Voltage-collapse model}

The four-dimensional state is
$[\delta_m,\omega,\delta,V]^{\mathsf T}$. The implemented Dobson--Chiang model
is~\cite{Dobson1989,Kong2021}
\begin{align}
  \dot{\delta}_m &= \omega,\\
  M\dot{\omega} &= -d_m\omega+P_m-E_m VY_m\sin(\delta_m-\delta),\\
  K_{qw}\dot{\delta} &= -K_{qv2}V^2-K_{qv}V+Q-Q_0-Q_1,\\
  TK_{qw}K_{pv}\dot V
  &=K_{pw}K_{qv2}V^2+(K_{pw}K_{qv}-K_{qw}K_{pv})V\nonumber\\
  &\quad+K_{qw}(P-P_0-P_1)-K_{pw}(Q-Q_0-Q_1),
  \label{eq:voltage}
\end{align}
where
\begin{align}
  P&=-E'_0 VY'_0\sin\delta+E_mVY_m\sin(\delta_m-\delta),\\
  Q&= E'_0 VY'_0\cos\delta+E_mVY_m\cos(\delta_m-\delta)
  -(Y'_0+Y_m)V^2.
  \label{eq:pq}
\end{align}
With $\theta_0=0$, the transformed source parameters are evaluated from
\begin{align}
  Y'_0&=Y_0\sqrt{1+(C/Y_0)^2-2(C/Y_0)\cos\theta_0},\\
  E'_0&=\frac{E_0}{\sqrt{1+(C/Y_0)^2-2(C/Y_0)\cos\theta_0}}.
\end{align}

\begin{table}[ht]
\centering
\caption{Voltage-model parameters used in every simulation.}
\label{tab:voltage-parameters}
\small
\begin{tabular}{@{}llllllll@{}}
\toprule
$K_{pw}$ & $K_{pv}$ & $K_{qw}$ & $K_{qv}$ & $K_{qv2}$ & $T$ & $P_m$ & $d_m$ \\
\midrule
0.4 & 0.3 & -0.03 & -2.8 & 2.1 & 8.5 & 1.0 & 0.05 \\
\addlinespace
$M$ & $E_m$ & $Y_0$ & $Y_m$ & $P_0$ & $Q_0$ & $P_1$ & $(C,E_0)$ \\
\midrule
0.01464 & 1.05 & 3.33 & 5.0 & 0.6 & 1.3 & 0.0 & $(3.5,1.0)$ \\
\bottomrule
\end{tabular}
\end{table}

The ODE is integrated by RK45 at $\Delta t=0.05$, with maximum internal step
$0.1$, relative tolerance $10^{-4}$, and absolute tolerance $10^{-6}$. Stable
training trajectories use
\begin{equation}
  \delta_m(0)\sim U(0.17,0.25),\quad
  \omega(0)\sim U(0,0.05),\quad
  \delta(0)\sim U(0.05,0.10),\quad
  V(0)\sim U(0.83,0.87).
\end{equation}
The first 500 time units (10,000 samples) are discarded. The next 3,000 time
units (60,000 samples) form the training record; an additional 200-time-unit
(4,000-sample) buffer is generated but is not included in the readout fit. The
conditioning statistics are measured over all 60,000 training samples. For the
displayed post-collapse ground truth, a delayed-collapse initial-condition
region is used:
$\delta_m(0)\sim U(0.17,0.42)$, $\omega(0)=0.025$,
$\delta(0)=0.10$, and $V(0)\sim U(0.70,0.87)$.

\subsection{Degree-frequency correlated Kuramoto network}

The target network obeys
\begin{equation}
  \dot\theta_i=\omega_i+K\sum_{j=1}^{N}A_{ij}
  \sin(\theta_j-\theta_i),\qquad
  \omega_i=\frac{k_i}{\langle k\rangle},
  \qquad i=1,\ldots,N.
  \label{eq:kuramoto}
\end{equation}
No additional $1/N$ normalization is applied to the coupling term. A fixed
Barabasi--Albert graph~\cite{Barabasi1999} is generated with $N=100$,
attachment parameter $m=3$, and network seed 0. It contains 291 edges, has mean
degree $5.82$, and has degree range 1--32. The degree-frequency correlation
produces explosive synchronization in this variant of the Kuramoto
model~\cite{Kuramoto1984,GomezGardenes2011}.

The ODE is integrated by RK45 with $\Delta t=0.05$ and relative and absolute
tolerances $10^{-7}$. For every training condition, the target is evolved for
185 time units. The first 70 time units (1,400 sampling intervals) are discarded
as target transient, leaving 115 time units (2,301 samples, including both
endpoints) for statistics and RC training. All five selected training
trajectories use target seed 26 and remain asynchronous. The final inter,
pre-synchronization trajectories displayed in Fig.~9 also use target seed 26.
The post query in Fig.~8 is assessed by a separate 20-window SCDT check, not
by a displayed seed-26 post trajectory. Selected-point GT checks use independent
phase seed 26, rather than the forward-continuation initialization of the GT rate curve.
The 10 highest-degree nodes form
$\mathcal H$ and the other 90
nodes form $\mathcal L$.

The target forward continuation spans $K\in[0.15,0.35]$ in steps of $0.01$;
the final rate figure displays its 12 GT points in $K\in[0.19,0.30]$.
The SCDT is evaluated at 18 uniformly spaced conditioning settings over the
same calibration-label interval. Each of 20
target realizations starts from independent random phases at $K=0.15$.
At each coupling, the target is integrated for 500 time units, and synchronization is
evaluated over the final 150 time units. The final phase state is then reused as
the initial condition at the next value on the increasing-$K$ grid. No
synchronized-branch initialization is imposed. The coupling value and
continuation state are not supplied as SCDT inputs.

\section{Training and test coordinates}

\begin{table}[H]
\centering
\caption{Food-chain coordinates. The SCDT receives the two statistics, not
$K$. Training statistics are direct measurements from the selected ODE
trajectories.}
\label{tab:food-coordinates}
\small
\begin{tabular}{@{}lrrrrl@{}}
\toprule
Role & $K$ & $\sigma_P$ & $H_{\mathrm{perm}}$ & Target seed & Coordinate source \\
\midrule
Train & 0.9740 & 0.121337 & 0.505563 & 8 & Direct target measurement \\
Train & 0.9780 & 0.123094 & 0.506485 & 5 & Direct target measurement \\
Train & 0.9800 & 0.125328 & 0.508786 & 1 & Direct target measurement \\
Train & 0.9820 & 0.126225 & 0.510812 & 0 & Direct target measurement \\
Train & 0.9880 & 0.128283 & 0.512399 & 8 & Direct target measurement \\
\rowcolor{bluegray} Inter & 0.9850 & 0.127254 & 0.511310 & -- & Conditioning-path query \\
\rowcolor{bluegray} Pre & 0.9970 & 0.133339 & 0.517649 & -- & Conditioning-path query \\
\rowcolor{bluegray} Post & 1.0001 & 0.134923 & 0.519299 & -- & Conditioning-path query \\
\bottomrule
\end{tabular}
\end{table}

\begin{table}[H]
\centering
\caption{Voltage-collapse coordinates. The SCDT receives the two statistics,
not $Q_1$.}
\label{tab:voltage-coordinates}
\small
\begin{tabular}{@{}lrrrl@{}}
\toprule
Role & $Q_1$ & $\sigma_V$ & $\rho_1(V)$ & Coordinate source \\
\midrule
Train & 2.9896800 & 0.0316339 & 0.9834461 & Direct target measurement \\
Train & 2.9897050 & 0.0317339 & 0.9835333 & Direct target measurement \\
Train & 2.9897300 & 0.0319752 & 0.9838649 & Direct target measurement \\
Train & 2.9897550 & 0.0322487 & 0.9841907 & Direct target measurement \\
Train & 2.9897800 & 0.0325905 & 0.9846133 & Direct target measurement \\
\rowcolor{bluegray} Inter & 2.9897475 & 0.0321755 & 0.9841014 & Direct $\sigma_V$; projected $\rho_1$ \\
\rowcolor{bluegray} Pre & 2.9897840 & 0.0326679 & 0.9847096 & Direct $\sigma_V$; projected $\rho_1$ \\
\rowcolor{bluegray} Post & 2.9898270 & 0.0330095 & 0.9851316 & Projected near-critical extrapolation query \\
\bottomrule
\end{tabular}
\end{table}

For Voltage, the five displayed training pairs are direct measurements from
the 60,000-sample records and exactly match the normalized statistics stored
with the trained reservoir. Statistic normalization retains the coordinatewise
minimum and maximum from the original five 10,000-sample training records.
These bounds are reused for the final 60,000-sample training records and all
test and rate queries, without refitting the normalization. Their fitted conditioning relation is
$\rho_1=1.2351899093\,\sigma_V+0.9443585599$. The inter, pre, and post
test queries, as well as every SCDT rate-sweep query, use this same relation.
The rate-sweep inputs are 35 uniformly spaced $\sigma_V$ values, and the
displayed rate axis is aligned by the Fig.~7(b) raw regression described
below. The measured-training fit $\sigma_V=9.7123093821\,Q_1-29.0051462885$
is used only to assign original query labels for selecting the training
warm-up record and conditioning schedule; $Q_1$ itself is used only for
target generation and axis alignment. Refitting the readout
on the 299,495 effective training pairs reproduced the stored readout weights
exactly; the training inputs and reservoir weights were not changed.

\begin{table}[H]
\centering
\caption{Kuramoto coordinates. Column order matches the conditioning vector
used by the code, $[S_{\mathrm{out}},S_{\mathrm{int}}]$.}
\label{tab:kura-coordinates}
\small
\begin{tabularx}{\textwidth}{@{}lrrrrY@{}}
\toprule
Role & $K$ & $S_{\mathrm{out}}$ & $S_{\mathrm{int}}$ & Target seed & Coordinate source \\
\midrule
Train & 0.2180 & 0.056196 & 0.183096 & 26 & Direct target measurement \\
Train & 0.2210 & 0.057167 & 0.203595 & 26 & Direct target measurement \\
Train & 0.2260 & 0.056519 & 0.208676 & 26 & Direct target measurement \\
Train & 0.2290 & 0.058910 & 0.216242 & 26 & Direct target measurement \\
Train & 0.2320 & 0.060497 & 0.227692 & 26 & Direct target measurement \\
\rowcolor{bluegray} Inter & 0.2300 & 0.059169325 & 0.218284892 & 26 & Training-regression query \\
\rowcolor{bluegray} Pre & 0.236042416 & 0.060820326 & 0.231408065 & 26 & Training-regression query \\
\rowcolor{bluegray} Post & 0.2400 & 0.061901678 & 0.240003312 & -- & Training-regression query \\
\bottomrule
\end{tabularx}
\end{table}

The Kuramoto training-coordinate fits are
$S_{\mathrm{out}}=0.273235303842\,K-0.003674794901$ and
$S_{\mathrm{int}}=7.948614050530\,S_{\mathrm{out}}-0.252029235654$.
Selected test pairs follow these fits and are the actual coordinates supplied
to the SCDT. During the pre search, $S_{\mathrm{int}}$ is scanned and
$S_{\mathrm{out}}$ is obtained by inverting the second fit; the first fit
then identifies the corresponding target-query $K$. These supplied test
coordinates are not newly measured raw statistics of the test trajectories.
The selected pre satisfies
$0.227692011<S_{\mathrm{int}}=0.231408065<0.231968081$, placing it
between the maximum training $S_{\mathrm{int}}$ and the displayed coarse-grid
SCDT crossing. Its original target-query $K=0.236042416$ exceeds the largest
training $K=0.232$. Its $S_{\mathrm{out}}=0.060820326$ also exceeds the
maximum raw training $S_{\mathrm{out}}=0.060497060$.
All 20 training-window SCDT reproductions remain asynchronous at this pre;
the independently integrated seed-26 GT trajectory is also asynchronous.
This is a single-realization GT check, not a 20-seed GT ensemble result.

For Food and Voltage inter and pre queries, the requested conditioning
coordinate is held fixed during warm-up and autonomous prediction. For
Kuramoto, warm-up pairs the selected training trajectory with its own training
conditioning coordinate; the requested query coordinate is applied from the
first autonomous step and held fixed thereafter. Food uses the last-training
conditioning for the first 500 autonomous steps and then ramps to the post
coordinate over 50 steps. Voltage and Kuramoto use a constant requested
conditioning coordinate throughout the autonomous rollout. Dynamic warm-up
samples are taken only from a selected training trajectory for all three
systems. Warm-up samples initialize the reservoir state but are excluded from
every autonomous evaluation window reported below.
\textcolor{red}{The Food rate sweep uses a single input schedule at every
conditioning setting: the requested coordinate is held constant during the
500-sample training-trajectory warm-up and all 1,000 autonomous samples.
There is no critical-parameter-dependent switch, autonomous hold, or ramp
in the rate sweep. The separately displayed post trajectory retains the
hold/ramp protocol described above.}

\begin{landscape}
\section{Sequence lengths and rollout protocol}

\centering
\captionof{table}{Complete sequence accounting. A ``step'' denotes one retained
sample interval after ODE resampling.}
\label{tab:sequence}
\small
\begin{tabularx}{\linewidth}{L{38mm} Y Y Y}
\toprule
Item & Food chain & Voltage collapse & Kuramoto network \\
\midrule
Sampling interval
& $\Delta t=1$
& $\Delta t=0.05$
& $\Delta t=0.05$ \\

Target transient discarded
& 2,000 steps = 2,000 time units
& 500 time units = 10,000 steps
& 70 time units = 1,400 intervals \\

Retained target record per training condition
& 6,000 steps
& 3,000 time units = 60,000 training samples, plus a 4,000-sample buffer
& 115 time units = 2,301 samples \\

Window used to calculate conditioning statistics
& First 2,000 samples of retained $P(t)$
& All 60,000 samples of the training-window $V(t)$
& All 2,301 retained phase samples \\

Training input/target pairs before washout
& 1,999 per condition
& 59,999 per condition
& 2,300 per condition \\

Training washout
& Updates $n\le100$ excluded
& First 100 updates excluded = 5 time units
& First 80 updates excluded = 4 time units \\

Readout pairs retained
& 1,898 per condition; 9,490 total
& 59,899 per condition; 299,495 total
& 2,220 per condition; 11,100 total \\

Prediction target seed buffer
& Selected inter/pre: final 500 samples from a training trajectory; rate: sampled 500-sample training windows
& 100 samples from the nearest selected training trajectory; endpoint 9,500 for inter/pre and 26,375 for the displayed post case; rate: 50 training-window endpoints per $Q_1$
& Displayed inter/pre: member 0 among 20 evenly spaced 60-sample windows from the nearest training trajectory; rate curve: all 20 windows \\

Samples actually fed during prediction warm-up
& 500 steps
& 100 steps = 5 time units
& 60 steps = 3 time units \\

Autonomous prediction horizon
& Selected trajectories: 2,000 steps; rate: 1,000 steps ($\Delta t=1$)
& Inter/pre and statistical reconstruction: 10,000 steps = 500 time units; post and rate: 2,000 steps = 100 time units
& 10,000 steps = 500 time units for both selected-point checks and the rate curve; first 2,000 steps shown in Fig.~9 \\

Post-query conditioning schedule
& Last-training conditioning held for 500 autonomous steps, followed by a 50-step linear ramp to the post coordinate
& Requested post coordinate fixed throughout the autonomous rollout
& Post coordinate fixed throughout the autonomous rollout \\

Displayed trajectory length
& 2,000 steps
& 1,000 steps for inter/pre; first 200 of 2,000 generated steps for the post comparison
& 2,000 steps \\
\bottomrule
\end{tabularx}
\end{landscape}

\begin{landscape}
\section{Reservoir hyperparameters and constraints}

\centering
\captionof{table}{Final reservoir settings. ``Density'' follows the exact mask semantics
of each implementation; this avoids interpreting the configuration field
\texttt{sparsity} in the opposite direction.}
\label{tab:hyperparameters}
\small
\begin{tabularx}{\linewidth}{L{41mm} Z Z Z}
\toprule
Parameter & Food chain & Voltage collapse & Kuramoto network \\
\midrule
Reservoir nodes $N_r$ & 1,000 & 900 & 2,000 \\
Recurrent-matrix construction
& Sparse random matrix with requested pre-symmetrization density $113/1000=0.113$; matrix is symmetrized
& Gaussian dense draw masked by $U>0.91$; approximately 9\% retained
& Gaussian draw with entries zeroed for $U>0.972323118$; approximately 97.23\% retained \\
Spectral radius & 1.5 & 0.5346534903 & 1.153870539 \\
Leakage $\alpha$ & 0.01 & 0.1426937444 & 0.09044233509 \\
Ridge coefficient $\beta$ & 0.00183011 & 0.08711372385 & 10.0 \\
Dynamic-input scaling & 5.0 & 2.566087561 & 0.04672287231 \\
Statistic-input scaling & $(0.1,0.1)$ after z-scoring & 0.125 in $W_{\mathrm{stat}}$ & 0.4542375643 in $W_{\mathrm{stat}}$ \\
Statistic offset & $(-1.631799943,0.0517568555)$ & None after min-max scaling & None after z-scoring \\
Reservoir seed & 11 & 42 & 23 \\
Output constraint
& Negative state components clipped to zero; rollout terminated at $P<0.01$ by filling the remaining output with zero
& Componentwise clipping to $[-0.10,-1.50,-0.30,0.01]$ and $[0.70,1.20,0.40,1.10]$
& Standardized output clipped to $\pm5.716762690$, then inverse-transformed and clipped to $r\in[0,1]$ \\
Prediction mode & Full-state closed loop & Full-state closed loop & Scalar $r$-only closed loop \\
\bottomrule
\end{tabularx}
\end{landscape}

The final Kuramoto checkpoint was obtained by exploratory searches over
five-point training sets and $N_r=2000$ reservoir settings, followed by
full-horizon, 20-window validation and ridge-coefficient refinement. Selection
used the frozen GT transition to favor an earlier SCDT crossing together with
asynchronous training/pre queries and a synchronized high-coordinate response.
These are exploratory model-selection results, not independent held-out
validation of the critical-point estimate. A separate $K=0.240$ post query
has subsequently been checked with all 20 training warm-up windows.
The standardized-output clipping ceiling maps to $r=0.550513278$ with the
stored signal normalizer; synchronization outcomes use this constrained output
and do not imply reproduction of the GT post-branch amplitude near $r=1$.
The inter and pre trajectories in Fig.~9 use the saved final checkpoint; the
rate curve in Fig.~10 uses the same checkpoint and training coordinates.
For the displayed inter and pre trajectories, the nearest training records
are $K=0.229$ and $K=0.232$, respectively. The first of the 20 training
windows is asynchronous in each case and is the window plotted in Fig.~9.

\section{Outcome criteria and rate curves}

\subsection{Selected-trajectory classification}

The qualitative selected-point checks use the following definitions:
\begin{itemize}
  \item \textbf{Food-chain survival:} a 2,000-step trajectory survives when
  the mean predator population over the final 100 steps is greater than 0.05.
  The first sample with $P<0.05$ is recorded as the event step.
  \item \textbf{Voltage survival:} a 2,000-step trajectory survives when every
  autonomous voltage sample is finite and $V>0.30$ throughout the rollout. The
  first sample with $V\le0.30$ is recorded as the event time.
  \item \textbf{Kuramoto synchronization:} the evaluation tail is the final
  30\% of a 10,000-step prediction, namely 3,000 samples or 150 time units. A
  trajectory is synchronized when at least 80\% of those samples satisfy
  $r\ge0.5$. Equivalently, at least 2,400 of the final 3,000 samples, corresponding
  to at least 120 time units in total, must pass the threshold. The passing
  samples need not be contiguous.
\end{itemize}

For the Kuramoto inter/pre checks, the autonomous prediction lasts 500 time
units after a 3-time-unit training-trajectory warm-up. Figure~9 displays only
its first 100 time units. Classification uses the final 3,000 autonomous
samples at $t=350.05,350.10,\ldots,500.00$, not the displayed first 100 or
the discarded 70-time-unit target transient. For the independent selected-point
GT, the matching evaluation samples are at absolute target times
$420.05,420.10,\ldots,570.00$ after the 70-time-unit transient.

For the final Kuramoto rate curve, the same threshold and occupancy fractions
are applied to a 10,000-sample autonomous prediction (500 time units). Its
evaluation tail is the final 3,000 samples (150 time units), and synchronization
requires at least 2,400 of those samples (120 time units in total) to satisfy
$r\geq0.5$. Warm-up samples are excluded from both evaluation windows.

\subsection{Rate-sweep grids and ensembles}

\begin{table}[ht]
\centering
\caption{Final rate-curve protocol and provenance.}
\label{tab:rate}
\small
\begin{tabularx}{\textwidth}{L{29mm} Y Y Y}
\toprule
Item & Food chain & Voltage collapse & Kuramoto network \\
\midrule
GT/calibration-parameter grid
& GT and SCDT calibration: $K\in[0.975,1.010]$, 22 equally spaced points each
& $Q_1\in[2.989620,2.989960]$, 35 points at $10^{-5}$ spacing
& GT: 12 points, $\Delta K=0.01$; SCDT calibration labels: 18 points, $\Delta K=0.11/17\approx0.006470588$, both in $[0.19,0.30]$ \\
Representative lower-axis statistic
& $\sigma_P$
& $\sigma_V$
& $S_{\mathrm{int}}$ \\
GT ensemble
& 20 fresh ODE integrations per final-grid condition, using fixed baseline initial-condition subsets
& 50 fresh ODE integrations per final-grid condition, using the original random-initial-condition family
& 20 forward-continuation ODE realizations per point \\
SCDT ensemble
& 20 closed-loop rollouts per point
& 50 closed-loop rollouts per point
& 20 closed-loop rollouts per point, initialized by distinct training-trajectory windows \\
GT/SCDT evaluation horizon
& 1,000 samples, $\Delta t=1$; 1,000 time units after initialization
& 2,000 samples, $\Delta t=0.05$; 100 time units
& 10,000 samples, $\Delta t=0.05$; 500 time units \\
Rate rule
& Final-500-step mean $P>0.05$
& All 2,000 predicted voltage samples finite and $\min V>0.30$
& Final 3,000 of 10,000 autonomous samples; at least 80\% with $r\geq0.5$ \\
Reported GT--SCDT MAE
& \textcolor{red}{21.1364} percentage points
& 17.7143 percentage points
& 0.03560606 in rate fraction, on the 12 GT coordinates \\
\bottomrule
\end{tabularx}
\end{table}

For the Kuramoto rate MAE, the 18-point SCDT curve is linearly interpolated
at the 12 directly evaluated GT coordinates on the common statistic axis.
The GT rates are not interpolated or recomputed for this comparison.
These are distinct conditioning grids, not different ensemble sizes.
For Food and Voltage, GT and SCDT are evaluated at matching abscissae on
the common statistical axis. The 22 Food supplied values are uniformly
spaced over $[0.119237869,0.149162935]$ with spacing $0.001425003$;
the 35 Voltage values cover $[0.030830874,0.034524236]$ with spacing
$0.000108628$. The same panel-(b) regression maps the corresponding GT
parameter grid onto these coordinates. MAEs compare the paired measured
rates directly, without interpolation. All 440 Food and 1,750 Voltage
SCDT member rollouts were regenerated with the frozen reservoirs and
training-only warm-up protocols; every evaluated point is retained.
\textcolor{red}{All 440 Food rollouts use the same constant-requested-coordinate
rule. The frozen reservoir, requested coordinates, training source records,
20 warm-up windows per setting, and survival criterion are unchanged.}

The final Food GT rate retains the 22 directly integrated conditions in
$K\in[0.975,1.010]$ from the previous 31-condition GT calculation;
it is not an interpolation of stored survival fractions. Initial
conditions are the fixed 20-member subset selected with seed 20260922 from the
original baseline at the nearest original $K$. No initial conditions are
filtered by their outcome in the new integrations. After a 500-time-unit
target transient, 1,000 samples at $\Delta t=1$ are evaluated. The survival
mean uses the final 500 samples. At each matched statistical coordinate,
the 20 SCDT runs use the same evaluation length, sampling interval, tail,
and threshold, excluding their
500-sample training-trajectory warm-up. All rate-query entropy values use the
same training $\sigma_P$--$H_{\mathrm{perm}}$ fit as the selected test queries.
This does not equate GT random initial conditions with SCDT warm-up states.

The GT evaluation samples are $t=500,501,\ldots,1499$, with the mean
calculated at $t=1000,\ldots,1499$. For SCDT, autonomous sample indices
$0,\ldots,999$ are evaluated and indices $500,\ldots,999$ enter the mean;
no teacher-forced samples enter either SCDT evaluation window.
The GT rates use the first 1,000 samples of the saved ODE records.
All 22 SCDT conditions were freshly evaluated over 1,000 autonomous samples;
the saved member trajectories reproduce every reported survival decision.
Selected inter/pre/post trajectory checks retain their separate
2,000-sample horizon and final-100-sample criterion; the displayed post
rollout also retains its distinct conditioning hold/ramp protocol.

For Voltage, the 50 directly integrated target trajectories at each of the
35 displayed $Q_1$ values are retained. Initial conditions retain the original baseline's
distribution and seed convention, using the nearest original baseline row
to identify the seed family, with no outcome-based resampling. Both GT and
SCDT use a 100-time-unit evaluation horizon at $\Delta t=0.05$. GT samples
are evaluated from $t=0$ through $99.95$, so early threshold crossings count;
no initial 25-time-unit GT interval is discarded. SCDT evaluation starts
after its 100-sample training-trajectory warm-up. A sampled $V\le0.30$
establishes non-survival even if integration later truncates. All 906
truncated GT runs contain such a sampled violation, leaving no unresolved
decisions. The frozen consistent-final SCDT is evaluated afresh at all
35 uniformly spaced input statistics, with 50 member rollouts per point;
their saved trajectories verify every reported member decision.

Kuramoto target
and SCDT rates use different but explicit
initialization protocols. The target uses 20 independent
forward continuations. At each supplied coordinate, the SCDT is reset and warmed up with 20
evenly spaced 60-sample windows from the nearest training trajectory; the
warm-up samples are excluded from the 10,000-step autonomous evaluation. These
windows are different reservoir initializations, not independent target-system
seeds or independent reservoir-weight realizations.

The red dashed lines in Figs.~4(a), 7(a), and 10(a) mark the respective
\emph{SCDT} 50\% crossings. For Food, the crossing is at the supplied
$\sigma_P=\textcolor{red}{0.134200402}$, whose Fig.~4(b) axis-aligned label is
$K=\textcolor{red}{0.992500000}$. The target 50\%
crossing is at actual $K=1.000000000$, mapped to $\sigma_P=0.140612916$
by the same Fig.~4(b) regression. This actual-parameter crossing equals
the physical reference $K_c=1.0000$ on the displayed GT grid.
The separate selected Food post query has supplied $\sigma_P=0.134923$,
beyond the SCDT rate crossing; its target is generated at $K=1.0001$ and
collapses under the selected-trajectory protocol. Its supplied coordinate
is not a direct measurement of the Fig.~4(b) regression at that target $K$.

\Needspace{11\baselineskip}
The final Voltage sweep uses the measured-training conditioning relation given above
at all 35 supplied statistic values. At each value, the SCDT is independently reset and
evaluated from 50 training-trajectory warm-up windows, using the 2,000-step
survival rule. The empirical SCDT 50\% crossing is
$\sigma_V=0.032734585$, aligned with $Q_1=2.989795250$ using Fig.~7(b).
The target 50\% crossing is at actual $Q_1=2.989825135$, mapped to
$\sigma_V=0.033059222$ by Fig.~7(b). The matched-statistic
GT--SCDT rate MAE is 17.7143 percentage points. Figures 5--6, the selected test queries, and this rate curve use the same
measured-training $\sigma_V$--$\rho_1$ relation for the second input
coordinate. The separate raw target measurements in Fig.~7(b) do not
serve as SCDT inputs.

\textcolor{red}{The training-parameter mapping identifies the training warm-up
record; it does not determine the displayed upper axis. For the Food rate
curve, original query labels do not select a conditioning schedule: the
requested coordinate is constant at every setting.} The final Food supplied interval corresponds to
original query labels $[0.969414215,1.027956992]$; the Voltage interval
corresponds to $[2.989605872,2.989986149]$. These original labels differ
from the displayed calibration intervals. The requested first statistic
and its training-fit second statistic are the actual SCDT inputs.

For Kuramoto, the validation-axis calibration is the raw-panel regression
$f_b(K)=1.057809579698\,K-0.020350615386$ from Fig.~10(b).
GT points are placed at $f_b(K_{\mathrm{GT}})$; SCDT points retain their
actual supplied $S_{\mathrm{int}}$; the upper axis displays $f_b^{-1}$.
The 18 SCDT input settings are uniformly spaced over
$S_{\mathrm{int}}\in[0.180633205,0.296992259]$, corresponding to the
calibration-label interval $[0.19,0.30]$ with step $0.11/17$.
The actual statistic spacing is $\Delta S_{\mathrm{int}}=0.006844650$.
The original
training-based query $K$ used to select a warm-up record is different from
this calibration label. It is obtained by inverting the two training fits,
and spans $[0.212663708,0.266239914]$ for these 18 settings.
The selected pre has calibration label $0.238000000$, while its independent
target trajectory is generated at $K=0.236042416$.

\subsection{Direct target-statistic checks in panels (b)}

The points in Figs.~4(b), 7(b), and 10(b) are \emph{direct target
measurements}, not the regression-produced conditioning values on the lower
axis of panels (a). Their dashed lines are least-squares regressions of the
additional target measurements displayed in the respective panels. These
raw-statistic regressions describe the local parameter--statistic trends
and are also the validation-axis mappings in all three corresponding
panels (a): GT abscissae are $f_b(p_{\mathrm{GT}})$, SCDT abscissae remain
the actual supplied statistic, and the upper axis is $f_b^{-1}$.
The raw checks are separate from the five training
trajectories and do not enter the reservoir.

For Fig.~4(b), 30 surviving target trajectories from seed 43 are shown in the
$K\in[0.974,0.990]$ neighborhood. Their predator standard
deviations use 2,000 target samples each. For Fig.~7(b), 40 surviving seed-43
target trajectories are shown in the $Q_1\in[2.989675,2.989784]$ neighborhood;
their voltage standard deviations use 10,000 samples each.

For Fig.~10(b), each of 40 conditions in $K\in[0.190,0.225]$ reports the mean
of 20 direct target measurements of $S_{\mathrm{int}}$. Each measurement uses
the final 150 time units of a 500-time-unit forward continuation. Only
conditions for which all 20 target realizations remained asynchronous were
eligible; the 40 displayed conditions span the precritical interval. Raw
statistics can differ from supplied regression coordinates because they are
measured on different realizations and finite windows. These restricted panels
show local trends, not global linearity across the full rate sweep.

\subsection{Critical values and statistical-axis alignment}

The upper physical-parameter axis in the rate figures is a validation aid.
System-specific mappings use exactly the raw-measurement least-squares
regressions displayed in panels (b). For Food,
$f_b(K)=0.855001902518\,K-0.714388986243$; for Voltage,
$f_b(Q_1)=10.862830916202\,Q_1-32.444905689842$.
The Kuramoto mapping is given above. These local fits are extrapolated
where the rate grid extends beyond the raw measurement intervals.
They do not replace the training stat-space fits that set the second
SCDT input coordinate, or alter the supplied first coordinate.
The upper-axis parameter labels are calibration labels and are never
supplied to the SCDT. Table~\ref{tab:critical} reports the empirical
target and SCDT 50\% crossings; the collapse-model physical references
are listed separately. We denote the mapped target crossing by
$s_c^{(\mathrm{GT})}$ and the
empirical SCDT crossing by $\widehat{s}_c^{(\mathrm{SCDT})}$; the two quantities
need not be numerically identical.

\begin{table}[H]
\centering
\caption{Critical values calculated from the rate files used in the main figures.}
\label{tab:critical}
\small
\begin{tabularx}{\textwidth}{L{31mm} Y Y Y}
\toprule
System & Physical reference and empirical target crossing & Empirical SCDT 50\% rate crossing & Interpretation \\
\midrule
Food chain
& Physical reference: $K_c=1.0000\mapsto\sigma_P=0.140613$; target 50\% crossing: $K=1.000000000\mapsto\sigma_P=0.140612916$
& Calibration $K=\textcolor{red}{0.992500000}$, $\sigma_P=\textcolor{red}{0.134200402}$
& Linear interpolation between adjacent SCDT survival-rate samples. \\
Voltage collapse
& Physical reference: $Q_{1,c}=2.9898256\mapsto\sigma_V=0.033064272$; target 50\% crossing: $Q_1=2.989825135\mapsto\sigma_V=0.033059222$
& Calibration $Q_1=2.989795250$, $\sigma_V=0.032734585$
& Linear interpolation between adjacent SCDT survival-rate samples. \\
Kuramoto network
& Target 50\% crossing: $K=0.240\mapsto S_{\mathrm{int}}=0.233523684$
& Calibration $K=0.238529412$, $S_{\mathrm{int}}=0.231968081$
& Crossings from the final 500-time-unit, 20-member rate curves. \\
\bottomrule
\end{tabularx}
\end{table}

These values should not be conflated. The mapped target crossing places
the empirical target outcome on the common statistical axis. Physical
references for the two collapse models have a separate definition.
The SCDT 50\% crossing is an empirical estimate obtained from its finite
ensemble and discrete rate grid.

For each rate curve, the reported crossing is obtained from the two adjacent
sampled points that bracket a rate of 0.5. If their coordinates and rates are
$(s_1,p_1)$ and $(s_2,p_2)$, the crossing is
\begin{equation}
  s_{50}=s_1+\frac{0.5-p_1}{p_2-p_1}(s_2-s_1).
\end{equation}
Thus the quoted SCDT critical coordinate is the linearly interpolated 50\%
crossing, whereas the target-aligned empirical value is the target 50\%
crossing mapped onto the same statistical axis.
For the same Kuramoto checkpoint, a separate 56-condition uniform sweep with
calibration spacing 0.002 gives $S_{\mathrm{int},50}=0.232465874$
(calibration $K=0.239$). Figure~10 and Table~\ref{tab:critical} use only
the displayed 18-condition SCDT grid and the unchanged 12-condition GT grid.
The SCDT crossing is a coarse interpolation,
not an additional simulation at the crossing coordinate. The new pre check
is reported separately and is not inserted into that 18-condition curve.
With the model, range, and selected pre fixed, candidate uniform SCDT grids
of 13--17 points were checked at their transition brackets; none placed the
displayed SCDT crossing both above the pre and below the frozen GT crossing.
The 18-point grid was the first to satisfy this display-order criterion and
was evaluated over its full range. This grid-selection result does not make
the interpolated crossing a resolution-independent critical value.

\section{Selected-point results and statistical self-consistency}

\subsection{Qualitative outcomes}

\begin{table}[ht]
\centering
\caption{Outcomes of the displayed selected trajectories. Where
reported, MAE is computed pointwise over the 2,000-step autonomous segment.}
\label{tab:outcomes}
\small
\begin{tabular}{@{}llccl@{}}
\toprule
System & Role & GT outcome & SCDT outcome & Trajectory result \\
\midrule
Food chain & Inter & Survived & Survived & MAE $=0.126414$ \\
& Pre & Survived & Survived & MAE $=0.156667$ \\
& Post & Collapsed & Collapsed & GT event step 411; SCDT event step 629 \\
\addlinespace
Voltage collapse & Inter & Survived & Survived & MAE $=0.033482$ \\
& Pre & Survived & Survived & MAE $=0.032724$ \\
& Post & Threshold crossed & Threshold crossed & GT event at 7.90; SCDT event at 8.85 time units \\
\addlinespace
Kuramoto network & Inter & Asynchronous & Asynchronous & MAE $=0.052154$ \\
& Pre & Asynchronous & Asynchronous & MAE $=0.109709$ \\
\bottomrule
\end{tabular}
\end{table}

The outcome classification, rather than pointwise agreement over a long
chaotic rollout, is the primary test for the transition experiments. All eight
displayed GT/SCDT trajectory pairs have the expected qualitative state. The
Kuramoto post coordinate is shown in the stat-space diagram and included in
an additional 20-window SCDT check, but no single post trajectory is displayed
in Fig.~9. The GT rate at $K=0.240$ is 0.50. The SCDT grid instead has rates
0.00 at calibration label $K=0.235294118$ and 1.00 at $K=0.241764706$;
$K=0.240$ is not an independently evaluated SCDT grid point. The additional
selected post check and these rate-curve settings should not be described
as a matched selected GT/SCDT post pair.

The displayed Voltage post condition is $Q_1=2.989827$, only
$1.4\times10^{-6}$ above the physical reference $Q_{1,c}=2.9898256$. The GT
uses collapse-basin initial-condition seed 28 and the SCDT uses the nearest
training trajectory ($Q_1=2.989780$) with a 100-sample warm-up ending at sample
26,375. The GT first satisfies $V\le0.30$ at autonomous step 158
($t=7.90$), and the SCDT first satisfies it at step 177 ($t=8.85$).
The SCDT remains below the threshold through the end of the displayed
200-step window. It can recover in a longer 2,000-step rollout, so the
displayed event is a threshold crossing rather than proof of permanent collapse.

\subsection{Supplied versus reproduced statistics}

For Food and Voltage, the statistics were remeasured from the closed-loop SCDT
trajectories at the inter and pre conditions. This test asks whether a supplied
conditioning coordinate is approximately reproduced by the generated dynamics;
the statistics are not imposed as hard constraints in Eq.~\eqref{eq:scdt-update}.

In the supplied-versus-reproduced standard-deviation panels (Figures 3(a) and
6(a) of the main manuscript), each candidate point denotes one conditioning
setting and one deterministic closed-loop rollout with the fixed final reservoir
seed. It is neither a different ODE seed nor an ensemble average. Food uses 40
noncollapsed conditioning settings and measures the reproduced standard
deviation over all 2,000 autonomous samples ($\Delta t=1$). The fixed plotting
window displays 36 of these 40 Food settings. Voltage evaluates 40
noncollapsed conditioning settings, of which 39 fall inside the fixed plotting
window, and measures each over all 10,000 autonomous samples
($\Delta t=0.05$, 500 time units). The 500-sample Food warm-up and the
100-sample Voltage warm-up are excluded from these evaluation windows, so no
teacher-forced samples contribute to the reproduced statistics.

\begin{table}[ht]
\centering
\caption{Self-consistency of the selected noncollapsed inter/pre trajectories.
Relative error is $(\text{reproduced}-\text{supplied})/\text{supplied}$.}
\label{tab:self-consistency}
\small
\begin{tabular}{@{}llrrrrrr@{}}
\toprule
System & Role & Supplied $\sigma$ & Reproduced $\sigma$ & Error & Supplied stat 2 & Reproduced stat 2 & Error \\
\midrule
Food & Inter & 0.127254 & 0.129448 & $+1.724\%$ & 0.511310 & 0.515384 & $+0.797\%$ \\
Food & Pre & 0.133339 & 0.132420 & $-0.689\%$ & 0.517649 & 0.508985 & $-1.674\%$ \\
Voltage & Inter & 0.0321755 & 0.0322780 & $+0.319\%$ & 0.984101 & 0.984219 & $+0.0119\%$ \\
Voltage & Pre & 0.0326679 & 0.0328875 & $+0.672\%$ & 0.984710 & 0.985125 & $+0.0421\%$ \\
\bottomrule
\end{tabular}
\end{table}

For Kuramoto, the dynamic SCDT output contains only $r(t)$ and does not contain
the 100 node phases. The pair $(S_{\mathrm{out}},S_{\mathrm{int}})$ therefore
cannot be remeasured from the SCDT output without extending the readout to the
node level.

\section{Reproducibility notes}

\begin{enumerate}
  \item Five pre-transition target conditions are used for every system. The
  physical parameter values label the numerical data-generation runs but are
  absent from the reservoir input.
  \item Statistics are calculated once per retained target segment and are held
  constant while that segment is processed by the reservoir.
  \item The readout is trained one step ahead and is then evaluated in closed
  loop. A low one-step training error is not used as a substitute for the
  long-horizon survival or synchronization test.
  \item Food and Voltage output clipping prevents numerical divergence from
  dominating the outcome test. The clipping bounds are listed in
  Table~\ref{tab:hyperparameters} and are part of the final protocol.
  \item Collapse trajectories are excluded when measuring the target-only
  representative-statistic mapping used for the upper/lower axis alignment.
  This filtering affects only the validation alignment, not the five SCDT
  training trajectories.
  \item The measured-training stat-space linearity coefficients are
  $R^2=0.96918$ for Food, $R^2=0.99912$ for Voltage, and $R^2=0.75565$ for
  Kuramoto.
\end{enumerate}

\section{References}

\begin{thebibliography}{99}

\bibitem{Jaeger2004}
H. Jaeger and H. Haas,
``Harnessing nonlinearity: Predicting chaotic systems and saving energy in
wireless communication,''
\emph{Science} \textbf{304}, 78--80 (2004).

\bibitem{Kong2021}
L.-W. Kong, H.-W. Fan, C. Grebogi, and Y.-C. Lai,
``Machine learning prediction of critical transition and system collapse,''
\emph{Physical Review Research} \textbf{3}, 013090 (2021).

\bibitem{McCann1994}
K. McCann and P. Yodzis,
``Nonlinear dynamics and population disappearances,''
\emph{The American Naturalist} \textbf{144}, 873--879 (1994).

\bibitem{Dobson1989}
I. Dobson and H.-D. Chiang,
``Towards a theory of voltage collapse in electric power systems,''
\emph{Systems \& Control Letters} \textbf{13}, 253--262 (1989).

\bibitem{Bandt2002}
C. Bandt and B. Pompe,
``Permutation entropy: A natural complexity measure for time series,''
\emph{Physical Review Letters} \textbf{88}, 174102 (2002).

\bibitem{Barabasi1999}
A.-L. Barabasi and R. Albert,
``Emergence of scaling in random networks,''
\emph{Science} \textbf{286}, 509--512 (1999).

\bibitem{GomezGardenes2011}
J. Gomez-Gardenes, S. Gomez, A. Arenas, and Y. Moreno,
``Explosive synchronization transitions in scale-free networks,''
\emph{Physical Review Letters} \textbf{106}, 128701 (2011).

\bibitem{Ziaeemehr2021}
A. Ziaeemehr and A. Valizadeh,
``Frequency-resolved functional connectivity: Role of delay and the strength
of connections,''
\emph{Frontiers in Neural Circuits} \textbf{15}, 608655 (2021).

\bibitem{Pereira2010}
T. Pereira,
``Hub synchronization in scale-free networks,''
\emph{Physical Review E} \textbf{82}, 036201 (2010).

\bibitem{Vlasov2017}
V. Vlasov and A. Bifone,
``Hub-driven remote synchronization in brain networks,''
\emph{Scientific Reports} \textbf{7}, 10403 (2017).

\bibitem{Kuramoto1984}
Y. Kuramoto,
\emph{Chemical Oscillations, Waves, and Turbulence}
(Springer, Berlin, 1984).

\end{thebibliography}

\end{document}
